\section{RT-2：从 Vision-Language Model 到 Vision-Language-Action Model}

PaLM-E 仍主要输出 language-level plan 或 answer。RT-2 更进一步：保留 VLM 的视觉语义能力，同时把机器人 action 离散化为 token，让 autoregressive decoder 直接预测低层动作。由此产生 Vision-Language-Action（VLA）模型这一接口。

\subsection{VLM prior 为什么有用}

VLM 通过 image-text pretraining 学习对象类别、属性、关系、常见用途和语言指令。机器人数据可能从未出现“把草莓放进正确颜色的碗”或“把灭绝动物移走”这类组合，但 VLM 已在互联网数据中学到 strawberry、color matching、extinct animal 等概念。RT-2 要测试这些 semantic prior 能否在 robot trajectory supervision 下进入 action prediction。

\lecturefigure{slide-24-rt2-transition.jpg}{从 VLM 语义世界进入 RT-2 机器人控制}{官方录像恢复幻灯片 V024；00:39:54}

\lecturefigure{slide-25-vlm-world-understanding.jpg}{VLM 同时编码 visual 与 semantic world knowledge}{官方录像恢复幻灯片 V025；00:40:42}

\begin{knowledgebox}{First use：VLA}
\term{Vision-Language-Action model（VLA）}接收视觉 observation 与语言 instruction，并输出机器人动作或动作 token。它与普通 VLM 的差异不只在最后一层 label：训练数据含闭环 robot trajectories，输出 token 具有控制语义，推理还受到控制频率、硬件和安全层约束。
\end{knowledgebox}

\subsection{把 action 表示成 token}

RT-1 已用离散 action bins 训练 Transformer policy；RT-2 把同一思想接到大型 VLM decoder。设连续动作第 $j$ 维范围为 $[l_j,u_j]$，用 $B$ 个 bins 离散化，则
$$
z_{t,j}=\operatorname{clip}\!\left(
\left\lfloor \frac{a_{t,j}-l_j}{u_j-l_j}B\right\rfloor,0,B-1
\right).
$$
其中 $a_{t,j}$ 是动作 $a_t$ 的第 $j$ 维，$z_{t,j}$ 是对应离散 token index，$l_j,u_j$ 是物理范围，$B$ 是 bin 数。推理时把 token 中心解码回连续动作，再交给 robot controller。

\lecturefigure{slide-26-vlm-as-robot-policy.jpg}{RT-1 与 RT-2：image + text 到 discretized actions}{官方录像恢复幻灯片 V026；00:45:15}

\lecturefigure{slide-27-action-tokenization.jpg}{Action tokenization：位置、旋转与 gripper 维度离散化}{官方录像恢复幻灯片 V027；00:45:57}

\begin{lstlisting}[caption={连续动作与 action token 的概念转换}]
def encode_action(action, low, high, bins):
    ratio = (action - low) / (high - low)
    token = floor(clip(ratio, 0.0, 1.0) * bins)
    return clip(token, 0, bins - 1)

def decode_action(token, low, high, bins):
    center = (token + 0.5) / bins
    return low + center * (high - low)
\end{lstlisting}

\begin{warningbox}{Action token 的三种误差}
第一是\textbf{quantization error}：bin 太少会损失精度；第二是\textbf{support error}：训练中未见的 dexterous motion 可能没有合适 token pattern；第三是\textbf{semantic collision}：若 action token 与文本 vocabulary 共用或映射不清，decoder 可能混淆输出语义。Tokenization 让模型可训练，不代表连续控制难题消失。
\end{warningbox}

\subsection{Autoregressive action 与 co-fine-tuning}

把一个时间步的动作写成 token sequence $z_{t,1:J}$ 后，模型预测
$$
p_\theta(z_{t,1:J}\mid I_t,g)
=\prod_{j=1}^{J}p_\theta(z_{t,j}\mid I_t,g,z_{t,<j}),
$$
其中 $I_t$ 是当前 image observation，$g$ 是语言 instruction，$J$ 是动作维度或 token 数，$z_{t,<j}$ 是此前已生成的 action tokens。Autoregressive factorization 借用了 language modeling machinery，但各维顺序和 decoding latency 会影响控制。

Co-fine-tuning 交替采样 web VQA/caption examples 与 robot trajectory examples。Web examples 维持 semantic prior；robot examples 教模型何时输出 action token。若只在 robot data 上训练，模型规模大却可能忘记通用语义；若只用 web data，则完全没有动作监督。

\lecturefigure{slide-28-rt2-training-data-models.jpg}{RT-2 training mixture：web-scale data + robot trajectories}{官方录像恢复幻灯片 V028；00:46:54}

\lecturefigure{slide-29-rt2-inference.jpg}{RT-2 inference：图像与指令经过 VLA 输出 action tokens}{官方录像恢复幻灯片 V029；00:48:54}

\teachervoice{讲者把 action representation 称为仍在探索的问题。Decimal、float-like string、extra tokens 都可以工作，但没有哪一种被宣告为普适最佳；若未来换新 embodiment 或新增 action dimension，旧 vocabulary 是否可迁移仍是开放问题。}

\begin{importantbox}{Co-fine-tuning 的因果链}
Web data 提供 semantic variation；robot data 把语义锚定到 action；共享参数允许两种 supervision 交互；ablation 再检查 improvement 是否真的来自 co-fine-tuning。只有这四步连起来，才能把“互联网知识进入控制”从叙事变成可检验机制。
\end{importantbox}

\subsection{本章小结}

RT-2 把 VLM 改造成 VLA：视觉和语言进入同一 decoder，低层动作离散为 token，web data 与 robot trajectories 共同训练。这个接口充分利用 autoregressive Transformer，但要承担 quantization、控制延迟、OOD action 与新 embodiment 迁移问题。

